% NOTA EDITORIAL: capítulo não incluído em main.tex nesta revisão. Conteúdo absorvido por chapters/08_discussao.tex (síntese) e appendices/apendice_f_custos.tex (agenda completa). Arquivo mantido apenas como registro histórico.

\chapter{IMPLICAÇÕES OPERACIONAIS E AGENDA DE CUSTOS}
\label{chap:implicacoes}

\section{Do perfil de cooperação preditiva à configuração de canais}

A motivação operacional deste trabalho parte de uma situação recorrente: sistemas de classificação dependem de múltiplas fontes de informação, mas a aquisição, a transmissão, o processamento e a manutenção dessas fontes têm custos distintos. Em tal contexto, a decisão relevante raramente é apenas identificar o subconjunto com menor perda em uma única execução. É necessário compreender quais configurações permanecem competitivas, quais relações de superioridade são estáveis e em que regime amostral os canais adicionais passam a justificar sua utilização.

Os resultados mostram que a correlação de ranking e a \WinnersAgreement{} são especialmente úteis para essa finalidade. A primeira descreve a preservação da organização global dos subconjuntos; a segunda informa diretamente se duas configurações mantêm a mesma relação de superioridade. Em conjunto, ambas permitem trabalhar com um grupo de alternativas eficientes, em vez de exigir a coincidência exata do primeiro colocado. Quando vários subconjuntos apresentam perdas próximas, uma decisão operacional pode preferir uma configuração ligeiramente inferior em desempenho, mas substancialmente mais barata, simples, robusta ou disponível.

Essa ideia pode ser formalizada, para cada tamanho amostral $n$, pelo conjunto de alternativas candidatas
\[
\mathcal{A}(n)=\left\{S_I:\overline{L}_{S_I,f}(n)\leq
\overline{L}_{\min,f}(n)+\varepsilon\right\},
\]
em que $\varepsilon$ representa uma tolerância definida pelo domínio. O \CoInfoSim{} não determina o valor dessa tolerância, mas fornece os rankings e as comparações par a par necessárias para construir $\mathcal{A}(n)$. Uma análise de custo pode então selecionar, entre as alternativas estatisticamente aceitáveis, aquela que melhor satisfaz as restrições operacionais.

A matriz de inversões $\Rmat$ e seus dois indicadores derivados, $\AR$ e $\SR$, acrescentam outra dimensão. Eles informam, respectivamente, se o modelo sintético recupera aproximadamente os mesmos pares que trocam de vencedor e se essas trocas ocorrem em regimes amostrais semelhantes dentro da grade observada. Essa informação é relevante para planejamento de coleta: um canal adicional pode não ser vantajoso em $n$ pequeno, mas tornar-se útil quando a operação acumula dados suficientes. O uso operacional de $\Rmat$, contudo, exige a cautela discutida nos capítulos anteriores: a matriz considera apenas a última inversão observada e não garante que a relação permanecerá inalterada além do maior $n$ avaliado.

\section{Custos não se reduzem ao número de canais}

A literatura de classificação sensível a custos distingue custos de erro, custos de teste, custos de aquisição e custos de decisão \cite{turney1995,turney2002}. Para o problema estudado aqui, a quantidade de canais é apenas uma aproximação muito rudimentar. Uma função de custo adequada deve ser definida para o domínio específico. Em termos abstratos, pode-se escrever
\[
C_{\mathcal{D}}(S_I,n;\theta_{\mathcal{D}}),
\]
em que $\mathcal{D}$ representa o domínio e $\theta_{\mathcal{D}}$ reúne parâmetros técnicos, financeiros e organizacionais. Essa notação não pressupõe uma forma aditiva nem proporcional.

Uma modelagem realista pode incluir:

\begin{itemize}
  \item investimento inicial em sensores, equipamentos ou infraestrutura;
  \item custo marginal por coleta, paciente, lote, ciclo ou unidade de tempo;
  \item calibração, manutenção, energia, transmissão e armazenamento;
  \item tempo de processamento e custo computacional;
  \item trabalho humano de inspeção, auditoria ou anotação;
  \item invasividade, risco, desconforto e indisponibilidade clínica;
  \item dependências entre canais produzidos pelo mesmo procedimento;
  \item custo do rótulo ou da referência usada para construir $Y$;
  \item custos de erro assimétricos e restrições regulatórias.
\end{itemize}

A presença de infraestrutura compartilhada impede assumir que o custo de um subconjunto seja a soma dos custos individuais. Um único equipamento pode produzir vários canais; nesse caso, retirar uma variável derivada não elimina o custo do procedimento. Inversamente, dois canais aparentemente semelhantes podem exigir sensores independentes, calibrações distintas e manutenção separada. Também não é correto assumir que o custo de $Y$ seja simplesmente proporcional a $n$: um rótulo pode depender de acompanhamento longitudinal, análise laboratorial em lote, auditoria especializada ou sensor de referência já instalado.

\begin{interpretationbox}[Alcance da análise de custos]
O presente trabalho não propõe uma função universal de custo nem demonstra que os subconjuntos identificados gerem economia. Sua contribuição é anterior: organizar o perfil de desempenho ao qual uma função de custo específica poderá ser acoplada.
\end{interpretationbox}

\section{Exemplos operacionais por domínio}

\subsection{Monitoramento ambiental}

No cenário Air Quality, os cinco canais PT08 representam respostas de sensores químicos, enquanto a concentração de benzeno de referência é utilizada para construir o alvo e não como preditor. Uma aplicação operacional pode envolver sensores de baixo custo distribuídos em vários pontos e um analisador certificado usado para calibração ou supervisão \cite{devito2008,UciAirQuality2008}. Nesse contexto, o custo pode depender de quantidade de estações, frequência de calibração, consumo energético, reposição de sensores, conectividade e disponibilidade do equipamento de referência.

Os resultados mostram que subconjuntos pequenos podem permanecer próximos do topo do ranking. Uma análise de custo poderia investigar se a perda adicional de uma configuração reduzida é aceitável diante da economia de aquisição ou manutenção. A alta preservação do ranking e da \WinnersAgreement{} em algumas condições sugere que classificadores treinados com amostras sintéticas produzidas por geradores parametrizados a partir dos dados reais podem ajudar a estudar esse espaço de alternativas sem depender exclusivamente de novas campanhas de coleta. A evidência, porém, não autoriza a retirada imediata de sensores: tal decisão exigiria validação temporal, análise de deriva da distribuição e estimativa financeira local.

\subsection{Ambientes construídos}

No Occupancy, os canais temperatura, umidade, luminosidade, concentração de CO$_2$ e razão de umidade podem estar associados a sistemas de automação predial e controle de HVAC \cite{candanedo2016,Weber2020}. O custo de uma configuração depende tanto dos sensores quanto da infraestrutura já presente no edifício. Temperatura e umidade podem ser fornecidas pelo mesmo dispositivo; luminosidade pode já estar disponível em sistemas de iluminação; CO$_2$ pode exigir calibração específica.

A geometria real do Occupancy mostrou-se distante da aproximação por uma única Gaussiana, mas os classificadores treinados com amostras dessa aproximação reproduziram ranking e relações de vencedores de forma competitiva em algumas configurações. Operacionalmente, isso ilustra que a utilidade de um modelo gerador não precisa ser julgada apenas pela fidelidade visual. Um modelo simples pode fornecer amostras suficientes para comparar alternativas de canais, mesmo quando não reproduz todos os agrupamentos da distribuição real. A escolha entre Gaussiana e GMM deve considerar qual componente da estrutura resultante do treinamento será usado na decisão.

\subsection{Aplicações clínicas}

No SUPPORT2, os canais combinam sinais vitais e exames laboratoriais. Esses grupos possuem estruturas de custo distintas: pressão arterial, frequência cardíaca, frequência respiratória e temperatura podem ser coletadas repetidamente à beira do leito; leucócitos, creatinina e sódio dependem de coleta, processamento laboratorial e tempo de resposta \cite{support1995,knaus1995}. Além do custo monetário, devem ser considerados invasividade, disponibilidade, urgência e consequência clínica do erro.

A elevada \WinnersAgreement{} do Random Forest indica que grande parte das comparações entre subconjuntos foi preservada nas condições sintéticas, mesmo com perdas absolutas altas. O resultado pode ser útil para estudar relações de eficiência entre combinações de sinais e exames. Contudo, o desempenho preditivo observado é insuficiente para sustentar decisões clínicas. A análise estrutural deve ser entendida como estudo metodológico das relações entre canais, não como validação de um sistema de prognóstico para uso assistencial.

\section{Uma formulação futura de decisão custo-eficiência}

Uma extensão natural consiste em combinar perda e custo em um problema multiobjetivo. Para um domínio $\mathcal{D}$, pode-se procurar o conjunto de configurações não dominadas segundo
\[
\left(
  \overline{L}_{S_I,f}(n),
  C_{\mathcal{D}}(S_I,n;\theta_{\mathcal{D}})
\right).
\]
Um subconjunto é dominado quando existe outro com custo não maior e perda não maior, com melhoria estrita em pelo menos uma dimensão. O resultado seria uma fronteira de Pareto de configurações custo-eficientes, abordagem relacionada à seleção de sensores e à otimização sob orçamento \cite{joshi2009,krause2008}.

As métricas de preservação poderiam então ser usadas como condição de confiança. Antes de otimizar sobre dados sintéticos, seria necessário verificar se a estrutura que emerge do treinamento na condição sintética reproduz suficientemente:

\begin{enumerate}
  \item a ordenação das alternativas relevantes;
  \item as comparações par a par próximas da fronteira eficiente;
  \item a existência e o regime amostral das inversões relevantes para planejar expansão amostral.
\end{enumerate}

Essa formulação também permite reconhecer que diferentes decisões exigem diferentes indicadores. Para escolher entre várias configurações de custo semelhante, $\rankrho$ e \WinnersAgreement{} podem ser prioritários. Para planejar quando instalar um canal adicional ou ampliar a coleta, $\AR$ e $\SR$ podem ser mais informativos. Não há necessidade de condensar todas essas propriedades em um índice único.

\section{Hipótese de projeção para regimes amostrais maiores}

A possibilidade mais ambiciosa do projeto é usar um gerador sintético validado no intervalo observado para explorar regimes amostrais maiores. O fluxo hipotético seria
\[
\begin{aligned}
\text{dados reais disponíveis}
&\longrightarrow \text{parametrização do gerador}
\longrightarrow \text{validação estrutural até }n_{\max}\\
&\longrightarrow \text{simulação em }n>n_{\max}
\longrightarrow \text{planejamento operacional}.
\end{aligned}
\]

Essa projeção não foi testada neste trabalho. A preservação nos tamanhos amostrais de 2 a 512 é condição necessária, mas não suficiente, para confiar em extrapolações. A estrutura pode mudar fora da faixa observada, especialmente quando o classificador entra em outro regime de variância, quando novos cruzamentos surgem ou quando o gerador não representa adequadamente as caudas e regiões de baixa densidade.

Uma validação futura deve ocultar deliberadamente parte da grade: parametrizar o gerador a partir dos dados disponíveis até um valor $n_0$, validar a estrutura resultante nesse intervalo, projetar os resultados para $n>n_0$ e comparar a projeção com simulações ou amostras reais posteriormente reveladas. Somente após esse teste será possível avaliar se a preservação observada pode apoiar planejamento de coleta, dimensionamento de canais ou decisões de investimento.

\section{Síntese das implicações}

O resultado operacional mais importante não é a identificação de um subconjunto universalmente ótimo. O trabalho fornece uma linguagem para analisar sistemas multicanais sob três perspectivas complementares: organização global das alternativas, relações par a par e transições ao longo de $n$. Essa estrutura pode ser integrada a modelos de custo, mas a integração deve respeitar as particularidades de cada domínio.

A agenda futura consiste, portanto, em transformar o \CoInfoSim{} de um simulador de preservação do perfil de cooperação preditiva em uma plataforma de análise custo-eficiência. A extensão deverá receber custos explícitos, identificar fronteiras não dominadas, quantificar incerteza econômica e testar projeções para regimes maiores. As conclusões atuais permanecem metodológicas: classificadores treinados com amostras sintéticas podem preservar componentes relevantes do perfil de cooperação preditiva, mas nenhuma economia ou decisão operacional específica foi demonstrada.
